\hypertarget{a-comprehensive-ai-learning-notev2.0}{%
\chapter{A Comprehensive AI Learning Note（v2.0）}\label{a-comprehensive-ai-learning-notev2.0}}

这是一份面向人工智能学习者的中文学习笔记暨讲义，内容覆盖深度学习、强化学习、大语言模型、大模型智能体、扩散模型、多模态生成、世界模型与具身智能等方向，覆盖面广于各类经典教材和市面上的绝大部分资料，且讲解详细、易懂。初版于2026年5月10日发布，当前为v2.0版本，共包含35章、168节。按正文（不含参考文献，以中文字符与英文/数字词计）统计，总字数约37.4万字。大家不必完整阅读，而是可以结合后面推荐入口中的目录和初学者学习路径，按需阅读自身感兴趣的内容。

本仓库由作者叶逸文的本地 Word 笔记转换而来。本仓库只保存可搜索、可阅读、便于协作纠错的 Markdown 版本及与之对应的 LaTeX 版本。中文为默认版本，英文版严格从中文版翻译，作者英文署名为 Yiwen Ye。

英文正文约18.7万词，按英文词计，排除参考文献、元数据、代码、图片、公式和链接目标。中文正文仍按中文字符与英文/数字词的既有口径统计。两种语言的可复现统计见 \texttt{tools/content\_statistics.py}。

\Needspace{5\baselineskip}
\hypertarget{ux5206ux4eabux76eeux7684}{%
\section{分享目的}\label{ux5206ux4eabux76eeux7684}}

\Needspace{5\baselineskip}
这份资料本来是个人的笔记，后经修改和补充，已成为可供学习者参考的讲义资料。将其公开的原因在于现在的教材虽然经典，但普遍存在两个问题：

\begin{enumerate}
\def\labelenumi{\arabic{enumi}.}
\item
  过于简洁。这会导致对于AI初学者而言理解难度较大，对有一定基础者而言又难以得到更有深度的理解。在我初次学习时，我通过和AI对话及通过不同资料交叉验证的方式学习，我也把这些思维过程均呈现在这份资料中。
\item
  没有与时俱进。AI时代发展日新月异，很多新领域几乎没有合适的学习资料。本讲义希望能填补这一空白。
\end{enumerate}

本讲义力求解决这些问题，尽可能详细，并覆盖当下AI发展的新方向、新领域。

\Needspace{5\baselineskip}
\hypertarget{ux5185ux5bb9ux8303ux56f4}{%
\section{内容范围}\label{ux5185ux5bb9ux8303ux56f4}}

\begin{itemize}
\tightlist
\item
  第1部分：深度学习（共9章、42节，约7.8万字）
\item
  第2部分：强化学习（共5章、25节，约5.4万字）
\item
  第3部分：大语言模型（共7章、47节，约9.8万字）
\item
  第4部分：大模型智能体（共4章、14节，约2.9万字）
\item
  第5部分：扩散模型与多模态生成（共4章、22节，约6.7万字）
\item
  第6部分：具身智能与世界模型（共6章、18节，约4.8万字）
\end{itemize}

\Needspace{5\baselineskip}
\hypertarget{ux7248ux672cux8bf4ux660e}{%
\section{版本说明}\label{ux7248ux672cux8bf4ux660e}}

\begin{itemize}
\item
  2026年5月10日发布 v1.0 版本。
\item
  2026年5月22日发布 v1.1 版本。

  更新内容：对部分图文进行了补充。
\item
  2026年9月2日发布 v2.0 版本（当前版本）。

  更新内容：对第2部分（强化学习）、第3部分（大语言模型）进行了微调，删除了一些非主流或学习价值不高的内容；对原第4部分、第5部分（涉及大模型智能体、扩散模型、多模态生成、具身智能、世界模型等）的内容进行了增删和修改，将其重构为当前的第4部分（大模型智能体）、第5部分（扩散模型与多模态生成）、第6部分（具身智能与世界模型），新增了业界的一些新方向、新趋势，删除了一些非主流或学习价值不高的内容，修改了部分文字表述。此外还修改了 \texttt{README.md}、初学者学习路径建议等。
\end{itemize}

\Needspace{5\baselineskip}
\hypertarget{ux63a8ux8350ux5165ux53e3}{%
\section{推荐入口}\label{ux63a8ux8350ux5165ux53e3}}

\begin{itemize}
\tightlist
\item
  \href{初学者学习路径.md}{初学者学习路径}
\item
  \href{目录.md}{完整目录}
\item
  \href{latex-project/README.md}{LaTeX 版本}
\item
  \href{DISCLAIMER.md}{免责声明}
\item
  \href{CONTRIBUTING.md}{贡献说明}
\end{itemize}

\Needspace{5\baselineskip}
\hypertarget{latex-ux7248ux672c}{%
\section{LaTeX 版本}\label{latex-ux7248ux672c}}

本仓库同时提供中英文 LaTeX 排版项目，位于 \href{latex-project/README.md}{latex-project}。该版本由公开 Markdown 内容整理生成，中文总书入口为 \texttt{latex-project/main.tex}，英文总书入口为 \texttt{latex-project/main-en.tex}；也可按节编号单独编译中文或英文 PDF。Markdown 版本仍作为主要阅读和维护版本，若两者存在差异，请以 Markdown 文档为准。

本地构建的 PDF 集中放在 \texttt{output/pdf/}：\texttt{complete-books/} 保存两本总书，\texttt{chinese-sections/} 和 \texttt{english-sections/} 各保存168个单节 PDF，\texttt{build-reports/} 保存构建报告，\texttt{PDF-INDEX.md} 提供总书与全部单节的可点击索引。两本总 PDF 同时保留在 \texttt{latex-project/build/}。PDF 和编译中间文件不上传到仓库；具体命令与输出位置规则见 \href{latex-project/README.md}{构建说明}。

\Needspace{5\baselineskip}
\hypertarget{ux8d21ux732eux4e0eux7ef4ux62a4ux8bf4ux660e}{%
\section{贡献与维护说明}\label{ux8d21ux732eux4e0eux7ef4ux62a4ux8bf4ux660e}}

欢迎通过 Issue 或 Pull Request 指出概念、公式、引用、图片转写和排版问题。作者目前为在校学生，如未及时查看或合并 PR，请谅解。感谢大家的支持！

\Needspace{5\baselineskip}
\hypertarget{ux8d44ux6599ux6765ux6e90ux4e0eux53efux9760ux6027}{%
\section{资料来源与可靠性}\label{ux8d44ux6599ux6765ux6e90ux4e0eux53efux9760ux6027}}

本讲义本来是个人学习笔记，包含作者整理、教材学习笔记、论文阅读笔记、业界动态整理，其中引用了学界和业界发布在Nature、AAAI、ICML、ICLR、NeurIPS、Arxiv等的部分论文、技术报告和其他发展动态，但可能存在漏标来源的情况。请以原论文、官方文档和权威教材为最终依据。

其中部分内容涉及个人解读，如有错误，敬请指正。

部分内容来自前沿研究论文，代表对应论文方法或作者理解，不应直接视为领域共识、工程最佳实践或业界标准范式。

\Needspace{5\baselineskip}
\hypertarget{ux8bb8ux53efux534fux8bae}{%
\section{许可协议}\label{ux8bb8ux53efux534fux8bae}}

本仓库采用 \href{https://creativecommons.org/licenses/by-nc-sa/4.0/}{Creative Commons Attribution-NonCommercial-ShareAlike 4.0 International} 协议。

可以在署名、非商用、相同方式共享的条件下复制、分享和改编本项目内容。完整法律文本以 \texttt{LICENSE} 和 Creative Commons 官方页面为准。

\clearpage
